%! Factoring Trinomials — Unit 1, Lesson 1.4 — Homework (Answer Key)
\documentclass[10pt]{article}
\usepackage{atda-article}
\usepackage{atda-key}   % pulls in -boxes; do NOT also load -boxes
\newcommand{\TallMath}[1]{$\displaystyle #1\rule[-1.4em]{0pt}{3.2em}$}
\setlength{\workrowsep}{8pt}

\begin{document}

\pageheader{Unit 1, Lesson 1.4}{Homework --- Answer Key}
% NAMESTRIP: no \namedateperiod here — mirrors the blank. Teacher-only prose goes
% in the lesson plan as \begin{teachernote}[Homework], never in a key.

\vspace{0.15in}

\boxguard[16]
\begin{notesbox}{Practice --- Factoring Trinomials}
\small
Show your work. Factor \textbf{completely} every time, and check by multiplying back. If no integer
pair works, write \emph{prime}.

\begin{enumerate}[leftmargin=1.6em, itemsep=6pt, topsep=4pt]

\item Factor: $x^2+8x+15$
\begin{work}
  3 \cdot 5 &= 15, \quad 3+5 = 8 \\
  x^2+8x+15 &= (x+3)(x+5)
\end{work}

\item Factor: $x^2-11x+24$
\begin{work}
  (-3)(-8) &= 24, \quad -3+(-8) = -11 \\
  x^2-11x+24 &= (x-3)(x-8)
\end{work}

\item Factor: $x^2+4x-21$
\begin{work}
  7(-3) &= -21, \quad 7+(-3) = 4 \\
  x^2+4x-21 &= (x+7)(x-3)
\end{work}

\item Factor: $x^2-5x-36$
\begin{work}
  (-9)(4) &= -36, \quad -9+4 = -5 \\
  x^2-5x-36 &= (x-9)(x+4)
\end{work}

\item Factor: $x^2+6x+11$
\begin{work}
  1 \cdot 11 &= 11, \quad 1+11 = 12 \\
  (-1)(-11) &= 11, \quad -1+(-11) = -12 \\
  x^2+6x+11 &\text{ is prime --- no integer pair sums to } 6.
\end{work}

\item Factor: $2x^2+11x+12$
\begin{work}
  ac &= 2 \cdot 12 = 24, \quad 3+8 = 11 \\
  2x^2+11x+12 &= 2x^2+3x+8x+12 \\
              &= x(2x+3) + 4(2x+3) \\
              &= (2x+3)(x+4)
\end{work}

\item Factor: $5x^2-17x+6$
\begin{work}
  ac &= 5 \cdot 6 = 30, \quad -2+(-15) = -17 \\
  5x^2-17x+6 &= 5x^2-2x-15x+6 \\
             &= x(5x-2) - 3(5x-2) \\
             &= (5x-2)(x-3)
\end{work}

\item Factor: $3x^2+5x-12$
\begin{work}
  ac &= 3(-12) = -36, \quad 9+(-4) = 5 \\
  3x^2+5x-12 &= 3x^2+9x-4x-12 \\
             &= 3x(x+3) - 4(x+3) \\
             &= (x+3)(3x-4)
\end{work}

\item Factor completely --- GCF first, then the trinomial: $4x^3+14x^2+6x$
\begin{work}
  4x^3+14x^2+6x &= 2x\left(2x^2+7x+3\right) \\
                &= 2x\left(2x^2+x+6x+3\right) \\
                &= 2x\left[x(2x+1) + 3(2x+1)\right] \\
                &= 2x(2x+1)(x+3)
\end{work}

\end{enumerate}
\end{notesbox}

\vspace{0.10in}

\boxguard[16]
\begin{scenariobox}[In Context --- The Garden Bed]{royal}
\small
The agriculture class is marking out a rectangular vegetable bed. The plan sheet gives only the
area:
\[ A(x) = 6x^2+13x+6 \quad \text{square feet,} \]
where $x$ is the number of feet the class extends the bed past the walkway.

\begin{enumerate}[leftmargin=1.6em, itemsep=6pt, topsep=4pt, start=10]

\item Factor $A(x)$ to recover the bed's two side lengths.
\begin{work}
  ac &= 6 \cdot 6 = 36, \quad 4+9 = 13 \\
  6x^2+13x+6 &= 6x^2+4x+9x+6 \\
             &= 2x(3x+2) + 3(3x+2) \\
             &= (3x+2)(2x+3)
\end{work}

\item The class settles on $x=4$. Give both side lengths and the area, in feet, with units --- then
check the area against the original polynomial.
\begin{work}
  3x+2 &= 14 \text{ ft}, \qquad 2x+3 = 11 \text{ ft} \\
  14 \cdot 11 &= 154 \text{ sq ft} \\
  6(16)+13(4)+6 &= 154 \text{ sq ft} \quad\checkmark
\end{work}

\item The class also has to buy edging for the border. In one sentence, say what the factored form
tells them that the expanded form does not.

\ansline{The side lengths --- $14$ ft and $11$ ft --- so the edging needed is $2(14)+2(11)=50$ ft.}

\end{enumerate}
\end{scenariobox}

\vspace{0.10in}

\boxguard[16]
\begin{scenariobox}[A First Look Ahead]{goldacc}
\small
\begin{enumerate}[leftmargin=1.6em, itemsep=6pt, topsep=4pt, start=13]

\item Two trinomials that today's method handles, and that Lesson 1.5 will give names to. Factor
each with the product-and-sum search you used above.
\begin{enumerate}[label=(\alph*), leftmargin=1.6em, itemsep=4pt, topsep=3pt]
  \item $x^2-9$ \quad --- first rewrite it as $x^2+0x-9$, then find two numbers with product $-9$
        and sum $0$.
\begin{work}
  3(-3) &= -9, \quad 3+(-3) = 0 \\
  x^2-9 &= (x+3)(x-3)
\end{work}
  \item $x^2+10x+25$
\begin{work}
  5 \cdot 5 &= 25, \quad 5+5 = 10 \\
  x^2+10x+25 &= (x+5)(x+5) = (x+5)^2
\end{work}
\end{enumerate}

\end{enumerate}
\end{scenariobox}

\vspace{0.10in}

\boxguard[18]
\begin{extensionbox}
\small
\textbf{Two variables, same search.} A landscaper writes a plot's area as $x^2+7xy+12y^2$ square
feet, where $x$ and $y$ are both measured in feet. Factor it into two binomials, then multiply your
answer back out. Why does the same product-and-sum search work when the constant term carries a
$y^2$?
\begin{work}
  3 \cdot 4 &= 12, \quad 3+4 = 7 \\
  x^2+7xy+12y^2 &= (x+3y)(x+4y) \\
  (x+3y)(x+4y) &= x^2+4xy+3xy+12y^2 \quad\checkmark
\end{work}
\ansline{The $y$ rides along untouched --- the search still only has to fix the numbers $3$ and $4$.}
\end{extensionbox}

\vspace{0.10in}

\begin{remindbox}
\small
\textbf{Next: Lesson 1.5 --- Special Forms} (\texttt{A2.EO.3b, d}). Item 13 already factored both of
them the long way. In 1.5 they get names --- the \textbf{difference of squares} and the
\textbf{perfect square trinomial} --- and patterns you can read straight off the page, plus the sum
and difference of cubes.
\end{remindbox}

\end{document}
